%======================================================================
\section{타당성을 위협하는 것}
\label{sec:threats}

\begin{yoji}
라벨은 층을 이루는 것과 같은 종류의 화살표이고 대부분 사용자 계정을 세요. 허락 층은 얇고, W1 동결에서는 주인 계약 하나가 그 대부분을 가져요. 주 분석은 등록 결과와 두 기간의 라벨을 안 뒤에 정했어요. 증거는 롤업 하나, 다섯 토큰, 라벨 기간 두 개에서 왔어요.
\end{yoji}

\subsection{구성 타당도}

\begin{itemize}
\item 라벨은 층을 이루는 것과 같은 종류의 화살표가 새로 도착하는 것이에요. 그래서 앞으로의 연결을 재고, 지금의 수를 따라가는 점수에 유리해요. 어느 쪽도 안전이나 손실의 결과가 아니고, “사용자가 기대는 계약”은 이 라벨로만 재요.
\item 라벨은 대부분 사용자 계정을 세는데, 어떤 점수도 그것을 읽지 않아요. 그래서 점수는 자기가 볼 수 없는 수요로 판단받아요.
\item 점의 규칙은 스마트 계정을 찾아내는 데 달려 있어요. 규칙은 Safe와 ERC-4337 계정을 찾아요. 목록의 코드 계열을 적용한 뒤에도, 계약으로 판별된 주소 115,084개는 \texttt{validateUserOp} 셀렉터는 있고 paymaster 셀렉터는 없는 코드를 가져요. 그리고 서비스가 사용자마다 또는 주문마다 만드는 계약도 따로 가려내지 않아요. 그런 계약이 점이면, 그래프에는 점의 규칙이 막으려던 종류의 계정이 들어 있는 거예요.
\item 종류는 자료가 끝난 뒤의 한 블록에서 코드로 읽어요.
\item 허락 화살표는 점수 기간의 마지막 허락을 동결 시점 주인의 잔액까지 센 것이지, 지금 유효한 허락이 아니에요(\ref{sec:model}절). 동결 시점에만 가진 잔액도 온전히 세어져요.
\end{itemize}

계약 사이의 허락 층은 얇아요.
\begin{itemize}
\item 2026년 6월 30일 동결에서 쌍이 115,006개이고 송금 쌍은 2,580,709개예요. 2026년 3월 31일 동결에서는 17,940개 대 2,269,073개예요.
\item 각 코호트의 5분의 1에서 4분의 1만 허락 쌍의 끝이에요. 그래서 허락 쪽의 모든 결과는 적은 계약에 기대요.
\item 2026년 6월 30일 동결에서는 주인 계약 하나가 115,006개 쌍 가운데 98,073개를 가져요. W1 점수와 같은 기간 점검은 이 계약 하나가 크게 모양을 정하는 층에서 계산했어요. 저자들은 그 계약이 무엇인지 알아보지 않았고, 그 계약 없이 분석을 되풀이하지도 않았어요.
\item 두 동결은 이 층에서 아주 달라서, W0과 W1은 허락 쪽에서 서로의 재현이 아니에요.
\end{itemize}
다른 쪽은 덮는 범위가 묶어요. 합집합 그래프 밖의 계약, 곧 스펜더 코호트마다 약 27\%는 만든 방식 때문에 저절로 0점의 동점이에요. 표~\ref{tab:union}은 그 계약들 없이 대비 두 개의 읽기가 달라진다는 것을 보여 줘요. 두 분석은 그래프, 무게 달기, 토큰, 점수 기간이 한꺼번에 달라서, 그 가운데 어느 하나를 훑은 것이 아니에요. 그리고 엔도스랭크와 AWP의 상한 있는 무게를 쓴 층 같은 다른 결합은 이 원고에서 평가하지 않아요.

\subsection{내적 타당도}

\begin{itemize}
\item 규칙 A는 이 연구가 로그를 하나도 추출하기 전에 정했어요. 규칙 B의 가설은 W0의 탐색적 대비에서, 그 허용 폭은 W0의 구간에서 나왔어요.
\item 등록은 독립된 등록처의 항목이 아니라 저자들 자신의 저장소에 있는 커밋이에요.
\item 주 분석은 모든 등록 결과를 알고, 두 기간의 라벨을 안 상태에서 정했어요. 그래서 그 선택이 어느 한 결과에 유리했을 수 있어요.
\end{itemize}
두 가지가 그 위험을 줄이지만 없애지는 못해요.
\begin{enumerate}
\item 상수는 W0 동결까지 12개월 동안 사건의 분위수에서 나와요. 그래서 라벨 기간의 사건은 상수에 들어가지 않아요.
\item 두 규칙에 대한 판정(둘 다 충족 못 함)이 등록 분석의 판정과 같고, A2, F1, F2의 읽기도 같아요. 두 분석 모두 A2는 실패, F1은 성립, F2는 실패예요.
\end{enumerate}
같은 것은 거기까지예요. 허락 걷기와의 대비인 A1과 F6a는 반대로, 계약 그래프에서 C-PR에 불리하게 읽혀요. A1은 주 분석에서 −0.025 [−0.045, −0.004], 등록 분석에서 +0.008 [0.003, 0.013]이고, F6a는 −0.043 [−0.061, −0.024]와 +0.010 [0.005, 0.014]예요.

주 분석을 먼저 보고하는 것도 결과 뒤에 한 선택이에요. 등록 분석은 규칙을 정할 때 시험하려던 분석으로서 규칙의 표와 그림에 남아 있어요. 이 원고는 규칙 B를 그 기준 두 개로 읽어요. 그 둘로는 규칙이 충족되지 않았음을 보일 수 있고, 충족됐음은 보일 수 없었어요. 그래프에 없는 계약을 0점으로 두는 규칙은 처리 과정의 선택이에요. S2는 그것이 어떤 읽기도 바꾸지 않는다는 것을, 표~\ref{tab:union}은 그것이 허락 걷기와의 대비에 무엇을 하는지를 보여 줘요.

\subsection{통계적 결론 타당도}

\begin{itemize}
\item 규칙에 묶인 구간은 계획에서 정한 대로 짝지은 다시 뽑기 400번에서 나와요.
\item 시드 50개를 더 써도 어느 분석에서든 A1, A2, F1, F2의 판정은 하나도 달라지지 않아요. 기준값에 가장 가까운 구간 끝, 곧 주 분석 A1의 위 끝도 모든 시드에서 0 아래에 머물러요(부록~\ref{app:explore}).
\item 규칙 옆에 보고한 많은 대비는 명목상이고, 시험의 가족으로 읽으면 안 돼요.
\item 계약 부트스트랩은 스펜더를 서로 독립으로 다뤄요. 하지만 새 주인 하나가 여러 스펜더에게 허락할 수 있어요.
\item 라벨은 듬성듬성해요. W1 스펜더 7,915개 가운데 1,491개만 새 허락 쌍을, 1,369개만 새 송금자를 얻었어요. 그래서 계수는 저마다 그 부분 집합들의 순서에 크게 기대요.
\end{itemize}

\subsection{외적 타당도}

증거는 롤업 하나, 다섯 토큰, 주인 셋 이상이 허락한 계약, 각 동결 전 1년으로 만든 점수(등록 분석에서는 2023년 10월부터의 모든 사건), 석 달짜리 라벨 기간 두 개, 한 블록에서 읽은 계정 종류에서 왔어요. 다른 체인이나 다른 토큰의 계약 그래프는 허락 층이 더 두꺼워서 다르게 움직일 수 있어요. 시험한 것은 균등 순간 이동을 쓰는 고정된 전체 무게뿐이에요. 점마다 다르거나 학습한 무게, 정상 벡터의 혼합, 층 사이 점의 복사본, 층마다 다른 감쇠 계수는 열린 문제로 남아요.
